\documentclass[12pt,a4paper]{article}

\usepackage[utf8]{inputenc}
\usepackage[spanish]{babel}
\usepackage{amsmath}
\usepackage{amssymb}
\usepackage{geometry}

\geometry{margin=2.5cm}

\title{Ejercicio 17: cantidad de operaciones y uso de memoria}
\author{}
\date{}

\begin{document}

\maketitle

Para analizar el costo del procedimiento, es necesario distinguir entre la generación de los puntos equiespaciados y la evaluación de los tres polinomios.

\section*{Generación de los puntos}

Para generar $100$ puntos equiespaciados entre $-1$ y $1$, se define primero la distancia entre dos puntos consecutivos:

\[
h = \frac{2}{100-1} = \frac{2}{99}.
\]

Luego, cada punto se calcula mediante:

\[
x_i = -1 + i h,
\qquad i=0,\ldots,99.
\]

Esto requiere:

\begin{itemize}
    \item una división para calcular $h$;
    \item $100$ multiplicaciones;
    \item $100$ sumas.
\end{itemize}

Por lo tanto, la cantidad total de operaciones aritméticas necesarias para generar los puntos es:

\[
\boxed{201\text{ operaciones}}.
\]

Si cada punto se almacena como un número de punto flotante de $64$ bits, cada uno ocupa $8$ bytes. Por lo tanto, los $100$ puntos ocupan:

\[
100 \cdot 8 = \boxed{800\text{ bytes}}.
\]

Este cálculo considera únicamente el espacio ocupado por los datos del arreglo y no incluye la sobrecarga interna introducida por las listas y los objetos de Python.

\section*{Evaluación de los polinomios por bloques}

El tamaño de bloque utilizado es $11$, debido a que el polinomio de mayor grado tiene grado $10$. Para procesar los $100$ puntos se necesitan:

\[
\left\lceil\frac{100}{11}\right\rceil = 10
\]

bloques.

Para cada bloque y cada polinomio se realizan las siguientes operaciones:

\begin{itemize}
    \item La construcción de una matriz de Vandermonde de tamaño $11\times11$ requiere $121$ potencias.
    \item La multiplicación de la matriz por el vector de coeficientes requiere $121$ multiplicaciones y $121$ sumas.
\end{itemize}

Por lo tanto, para cada bloque y cada polinomio se realizan:

\[
121+121+121=363
\]

operaciones.

Como se procesan $10$ bloques para cada uno de los $3$ polinomios, la cantidad total de operaciones para evaluarlos es:

\[
363\cdot10\cdot3
=
\boxed{10890\text{ operaciones}}.
\]

Al sumar las $201$ operaciones correspondientes a la generación de los puntos, se obtiene:

\[
10890+201
=
\boxed{11091\text{ operaciones}}.
\]

En este conteo, cada potencia se considera como una única operación. Si cada potencia se descompusiera en multiplicaciones elementales, la cantidad de operaciones sería mayor.

\section*{Espacio en memoria}

Considerando únicamente el espacio ocupado por los arreglos numéricos principales:

\begin{itemize}
    \item Los $100$ puntos ocupan:
    \[
    100\cdot8=800\text{ bytes}.
    \]

    \item Los resultados de los tres polinomios ocupan:
    \[
    3\cdot100\cdot8=2400\text{ bytes}.
    \]

    \item La matriz de Vandermonde de $11\times11$ ocupa:
    \[
    11\cdot11\cdot8=968\text{ bytes}.
    \]

    \item La matriz traspuesta ocupa otros:
    \[
    968\text{ bytes}.
    \]

    \item El vector de coeficientes y el vector resultante del bloque ocupan:
    \[
    2\cdot11\cdot8=176\text{ bytes}.
    \]
\end{itemize}

En consecuencia, el espacio aproximado ocupado por los arreglos principales es:

\[
800+2400+968+968+176
=
\boxed{5312\text{ bytes}}.
\]

Esto equivale aproximadamente a:

\[
\boxed{5.2\text{ KiB}}.
\]

Este valor no incluye la memoria adicional utilizada por las listas de Python, los objetos internos ni las estructuras creadas para realizar el gráfico.

\section*{Crecimiento con respecto a \texorpdfstring{$n$}{n}}

Si se construyera una única matriz de Vandermonde de tamaño $n\times n$, tanto la cantidad de operaciones como el espacio en memoria crecerían cuadráticamente:

\[
\boxed{\text{Operaciones}=\Theta(n^2)}
\]

y

\[
\boxed{\text{Memoria}=\Theta(n^2)}.
\]

En cambio, al trabajar con bloques de tamaño fijo $11$, la cantidad de bloques es proporcional a $n$. Por lo tanto, la cantidad de operaciones crece linealmente:

\[
\boxed{\text{Operaciones}=\Theta(n)}.
\]

El espacio auxiliar ocupado por la matriz de Vandermonde y su traspuesta permanece constante, ya que ambas tienen un tamaño máximo de $11\times11$. Sin embargo, almacenar todos los puntos y todos los resultados requiere una cantidad de memoria proporcional a $n$:

\[
\boxed{\text{Memoria total}=\Theta(n)}.
\]

\section*{Modificación para mejorar la eficiencia}

Para hacer el cálculo más eficiente, se podría modificar el orden de los ciclos. La matriz de Vandermonde depende únicamente de los puntos del bloque y no de los coeficientes del polinomio.

Por lo tanto, para cada bloque se puede calcular una sola vez la matriz de Vandermonde y su traspuesta:

\[
V=\operatorname{matrizVandermonde}(\text{bloque}),
\qquad
V^T=\operatorname{traspuesta}(V).
\]

Luego, la misma matriz traspuesta puede reutilizarse para evaluar los tres polinomios, cambiando solamente el vector de coeficientes.

De esta forma, se evita construir y transponer tres veces la misma matriz para cada bloque. La complejidad asintótica continúa siendo lineal con respecto a $n$, pero se reduce la cantidad concreta de operaciones realizadas y también la creación innecesaria de matrices temporales.

\end{document}